\subsection{Mechanical Waves} 

%\DTMnewdatestyle{dateupdated}{〈definition 〉

Answers to M\Alph{subsection} problems are on page \pageref{waves_prob_answers}.  \textbf{\textcolor{red}{Updated \DTMnow.}}

%--------------------------------------------------------------------------------------------------------------
\begin{Exercise}
A wave travels with speed of 200 m/s.  Its wave number is 1.5 rad/m.  What are (a) its wavelength and (b) its frequency?
\end{Exercise}
\begin{Answer}
(a) 4.2 m (b) 48 Hz
\end{Answer}


%--------------------------------------------------------------------------------------------------------------
\begin{Exercise}
The displacement of a wave traveling in the positive $x$ direction is
$$D(y,t)=(3.5\text{ cm})  \sin\Big[ \left( 2.7 \text{ m}^{-1} \right) x -  (124\text{ s}^{-1})t \Big] .$$
What are its (a) frequency, (b) wavelength, and (c) the speed of this wave?
\end{Exercise}
\begin{Answer}
(a) 20 Hz (b) 2.34 m (c) 45 m/s
\end{Answer}


%--------------------------------------------------------------------------------------------------------------
\begin{Exercise}
Consider a sinusoidal wave traveling on a rope.  The oscillator that makes the wave completes 40~cycles in 30~seconds.  A single crest of the wave travels 425~cm along the rope in 10~seconds.  What is the wavelength of the wave?  
\end{Exercise}
\begin{Answer}
0.319 m
\end{Answer}


%--------------------------------------------------------------------------------------------------------------
\begin{Exercise}
A string is driven at a frequency of 5~Hz to create a traveling wave moving to the right at speed $v=20.0$~m/s, with amplitude $A=12.0$~cm.  
Find (a)~the angular frequency and (b)~the wave number for this wave.  
(c)~Write a mathematical expression for the wave function.  
Calculate (d)~the maximum transverse speed and (e)~the maximum transverse acceleration of an element of the string. \textit{Hint for (d) and (e): if $y(t)$ is the position of a particle, then $y'(t)$ and $y''(t)$---that is, the first and second time derivatives---give the velocity and acceleration of the particle, respectively.}
\end{Exercise}
\begin{Answer}
(a) 31.4 rad/s (b) 1.57 rad/m (c) $y=(12\text{ cm}) \sin \left[ (1.57\text{ m}^{-1})x - (31.4\text{ s}^{-1})t \right] $ (d) 3.77 m/s 
(e) 118 m/s$^2$.
\end{Answer}


%--------------------------------------------------------------------------------------------------------------
\begin{Exercise}
A driver is traveling north on a highway at a speed $v=25$~m/s.  
A police car, traveling south at 40~m/s approaches with its siren blaring at a frequency of 250~Hz.  
(a)~What frequency does the driver hear as the police car approaches? 
(b)~What frequency does the driver hear after the police car passes? 
(c)~Repeat parts (a) and (b) for the case when the police car is behind the driver and travels north.  
\end{Exercise}
\begin{Answer}
(a) 304 Hz (b) 208 Hz (c) 262 Hz and 240 Hz
\end{Answer}


%--------------------------------------------------------------------------------------------------------------
\begin{Exercise}
Read the following documents on the ``Course Information'' Canvas Module:
\begin{itemize}[nosep]
\item The Syllabus
\item ``On Working Together Honestly''
\item Homework Guidelines
\item Why No Equation Sheets
\end{itemize}
(a)~If you could change any aspect of these course policies (or anything else about the course, I guess), what would you change?  
You must make at least one suggestion for a change; no fair just answering that you're fine with whatever's currently there.  
(b)~How serious are you about this request?  (Are you just suggesting a change because it was required, or because you genuinely think it would improve the course?) 
\end{Exercise}


%--------------------------------------------------------------------------------------------------------------
\begin{Exercise}
In class, I showed that applying $F=ma$ to a small element of a piece of string with mass density $\mu$, under tension $F_T$, leads to the ``wave equation'' in one dimension:
\[ \frac{\partial^2y}{\partial x^2} = \frac{\mu}{F_T}\frac{\partial^2y}{\partial t^2} \]
(a)~Show, by directly substituting into the equation above, that 
$y(x,t)=A\sin(kx+\omega t)$
satisfies the wave equation above, provided that 
$\omega^2/k^2=F_T/\mu$.
(This is almost exactly what we did in class, except that the ``$+$'' sign indicates that this wave is moving to the left instead of the right.) 
(b)~Is the equation for a standing wave, $y(x,t)=2A \sin(kx)cos(\omega t)$, a solution to the wave equation as well?  
(c)~Rewrite the ``wave equation'' above using the speed $v$ of the wave in place of the constants $\mu$ and $F_T$.  
\end{Exercise}


%--------------------------------------------------------------------------------------------------------------
\begin{Exercise}
The linear wave equation for a transverse wave in a solid material (such as the ``s-waves'' of an earthquake) is given by
\[ G\frac{\partial^2y}{\partial x^2} = \rho \frac{\partial^2y}{\partial t^2}, \]
where $G$ is a quantity called the ``shear modulus'' of the material (a measure of stiffness, or ``springiness,'' basically), and $\rho$ is the mass density (per volume) of the material.  Much of the subsurface bedrock in western Virginia is granite, which has a shear modulus of $G=24.5 \times 10^9$~N/m$^2$ and a density of 2700~kg/m$^3$.  
Find the speed of a transverse seismic wave in granite.  
\end{Exercise}
\begin{Answer}
Roughly 3000 m/s
\end{Answer}


%--------------------------------------------------------------------------------------------------------------
\begin{Exercise}
On a string with a tension of 6~newtons, transverse waves travel with a speed of 20~m/s.  What tension is needed for a transverse wave to have a speed of 30~m/s on the same string? 
\end{Exercise}
\begin{Answer}
13.5 N
\end{Answer}

%--------------------------------------------------------------------------------------------------------------
\begin{Exercise}
Two transverse sinusoidal waves described by the wave functions
\[y_1=3\text{~cm} \sin\Big[ \left( \pi \text{~m}^{-1}\right) x + \left( 0.6\pi \text{~s}^{-1} \right) t\Big] 
\quad \text{and} \quad
  y_1=3\text{~cm} \sin\Big[ \left( \pi \text{~m}^{-1}\right) x - \left( 0.6\pi \text{~s}^{-1} \right) t\Big] 
\]
combine to form a standing wave.  
Determine the maximum transverse position of an element of the medium at (a)~$x=0.25$~m, (b)~$x=0.5$~m, and (c)~$x=1.5$~m.  
(d)~Find the three smallest values of $x$ corresponding to antinodes.  
\end{Exercise}
\begin{Answer}
(a) 4.24 cm (b) 6 cm (c) 6 cm (d) 0.5 m, 1.5 m, 2.5 m
\end{Answer}


%--------------------------------------------------------------------------------------------------------------
\begin{Exercise}
In class I drew pictures on the whiteboard of the three lowest frequency standing waves in a pipe of length $L$ with one open end and one closed end.  
By looking at the pictures, I came up with a general formula for the wavelengths 
($\lambda = 4 L / n,$ where $n=1,3,5,...$) 
and a corresponding formula for the frequencies: $f=nv/4L$.  
(a)~Reproduce the drawings and these derivations.  
(b)~Do similar drawings and derivations for a pipe that's closed at both ends.  
(c)~Now do it for a pipe that's open at both ends.  
(You can check your work with Knight on page 465, but try to do them without looking at your book or your notes.  You may have to on a quiz!)
\end{Exercise}
\begin{Answer}
(b) $f=nv/2L,$ where $n=1,2,3,...$ which is the same as for a string 
(c) also $f=nv/2L,$ where $n=1,2,3,...$
\end{Answer}

%--------------------------------------------------------------------------------------------------------------
\begin{Exercise}
Students using the ``Economy Resonance Tube'' (open at one end, closed at the other, length 80~cm and radius 7.5~cm) measure resonances at 92~Hz, 297~Hz, 499~Hz, and 706~Hz.  
Assume these correspond to the odd integers $n=1,$ 3, 5, and 7, respectively.  
(a)~Use the ``LINEST'' function in Excel to find the slope of the graph of $f$ vs. $n$ (see Appendix \ref{excel} in your lab manual), and use that value to determine the speed of sound in air in the room, with uncertainty.  Be sure to use the radius to calculate the effective length $L_\text{eff}$, as described in Lab~\ref{resonance_in_tubes}. 
(b) To get the correct answer, the third argument of the LINEST function should be ``0'' instead of ``1''.  Why? \end{Exercise}
\begin{Answer}
$v=338.804 \pm 2.184$~m/s, which rounds properly to $v=339 \pm 2$~m/s
\end{Answer}

%--------------------------------------------------------------------------------------------------------------
\begin{Exercise}
On a piano, some keys cause the hammer to strike two or three strings at the same time, all tuned to the same note.  The note at 110~Hz has two strings, for instance.  If one of those two strings slips from its normal tension of 600~N to 540~N, what beat frequency is heard when that key is pressed? 
\end{Exercise}
\begin{Answer}
5.64 beats per second
\end{Answer}




\vspace{0.7in}

\textbf{Answers to Selected {\thesubsection} Problems:}
\label{waves_prob_answers}
%\shipoutExercise
\shipoutAnswer

\cleardoublepage
